\documentclass[11pt]{article}
\usepackage{amsmath,amssymb,amsthm}
\usepackage{geometry}
\usepackage{enumitem}
\usepackage{booktabs}
\usepackage{graphicx}
\geometry{margin=1in}

\begin{document}

\begin{center}
{\Large STCS 6701: Probabilistic Machine Learning\\
Homework 3}

\vspace{1em}

Due: Dec 5 at 11:59 pm ET
\end{center}

\section*{Abstract}

This project develops and evaluates a probabilistic latent factor model for urban property valuation under pervasive missingness and heavy-tailed noise. Using a large panel of New York City property records (roughly $10^5$ parcels) with assessment variables, building characteristics, and partial transaction data, I treat observations of sale prices as a noisy, mostly missing target and learn a low-dimensional latent representation that jointly explains both the feature distribution and prices. Concretely, I implement a semi-supervised variant of the Missing-data Importance-Weighted Autoencoder (MIWAE) with a mixture of Student-$t$ priors over the latent variables to better accommodate multimodality and extreme residuals. The model performs amortized inference for the latent variables, and produces both point predictions and predictive uncertainty for sale prices.

Quantitatively, I compare the generative model against strong baselines such as gradient-boosted trees using held-out predictive log-likelihood, RMSE, MAE, and median absolute percentage error (MAPE). I also assess calibration of the predictive variance and examine how much signal in prices is captured by the learned factors versus left in residuals. Qualitatively, I visualize the latent variables colored by building class and sale price to interpret what structure the model has discovered. The broader aim is to understand when a probabilistic latent factor model is competitive with standard supervised approaches for noisy, heavy-tailed real-estate data, and how its uncertainty estimates and latent structure could be used to flag unreliable valuations or systematic gaps in market coverage. Accurate, uncertainty-aware valuations are critical for equitable property taxation and urban planning because they help identify systematic assessment errors that may disproportionately affect specific neighborhoods.

\section*{Introduction}

\subsection*{(i) Problem and Motivation}

I implement a semi-supervised, probabilistic latent factor model to analyze a large panel of New York City property records. The primary goal is to learn a low-dimensional latent representation of parcels that captures spatial and physical heterogeneity across assessed values, land characteristics, and building attributes, and to use this representation to model log-transformed sale prices, including when prices are missing or partially observed.

Accurate property valuation is fundamental to municipal finance, real estate transactions, and equitable taxation. However, traditional automated valuation models (AVMs) often struggle with the heavy-tailed noise and pervasive missingness characteristic of urban data, leading to systematic errors that can disproportionately affect specific neighborhoods or property types. By explicitly modeling these factors, we aim to provide a more robust tool for assessors and policymakers.

The dataset is multivariate, noisy, and heavily skewed in both covariates and outcomes. A latent variable model with non-linear likelihoods is appropriate: it exploits shared structure across covariates, explicitly models missingness, and provides both point predictions and uncertainty estimates for log-transformed sale prices. I use a semi-supervised variant of a Missing-data Importance-Weighted Autoencoder (MIWAE) with a Student-\emph{t} mixture prior over the latent variables, trained on approximately $1.2\times 10^5$ NYC parcels.

\subsection*{(ii) Model}

Let $i = 1,\dots, n$ index properties. For each property we observe a covariate vector $x_i \in \mathbb{R}^D$ (parcel and building characteristics) and, for a subset, a scalar outcome $y_i = \log(\text{sale\_price}_i)$. The model introduces a latent factor $z_i \in \mathbb{R}^K$ with $K = 3$.

\paragraph{Prior on latent variables.} I place a $K$-component Student-\emph{t} mixture prior on $z_i$:
\[
c_i \sim \text{Categorical}(\pi), \qquad
z_i \mid c_i = k \sim t_\nu\big(\mu_k, \Sigma_k\big),
\]
where $\pi \in \Delta^{K-1}$ are mixture weights, $\nu > 0$ is the degrees-of-freedom parameter, and $(\mu_k,\Sigma_k)$ are component means and covariances. In the implementation, the prior is parameterized in terms of unconstrained logits and Cholesky factors. This parameterization induces a mixture of heavy-tailed components that can capture multi-modal, robust structure in latent space.

\paragraph{Likelihood for covariates.} Conditional on $z_i$, the observed covariates $x_i$ are generated by a diagonal-covariance Gaussian decoder:
\[
x_i \mid z_i \sim \mathcal{N}\big(\mu_x(z_i), \operatorname{diag}(\sigma^2_x(z_i))\big),
\]
where $\mu_x(\cdot)$ and $\log \sigma_x^2(\cdot)$ are outputs of a neural decoder network. In practice, each feature dimension is scaled to have approximately zero mean and unit variance before training, so the decoder operates in standardized units.

\paragraph{Likelihood for log-transformed sale price.} For the log-transformed sale price $y_i$, I add a supervised Gaussian head on top of $z_i$:
\[
y_i \mid z_i \sim \mathcal{N}\big(\mu_y(z_i), \sigma_y^2(z_i)\big),
\]
with $\mu_y(\cdot)$ and $\log \sigma_y^2(\cdot)$ given by a small prediction network that takes $z_i$ as input. This defines a conditional generative model for $y_i$ given $z_i$; the latent variables thus jointly explain both covariates and prices.

\paragraph{Joint distribution.} The full joint distribution of observed and latent variables for one property is
\[
p(x_i, y_i, z_i, c_i)
= p(c_i)\, p(z_i \mid c_i)\, p(x_i \mid z_i)\, p(y_i \mid z_i)^{\mathbb{I}\{y_i \text{ observed}\}},
\]
and across properties I assume conditional independence given the global parameters:
\[
p(X, Y, Z, C)
= \prod_{i=1}^n p(x_i, y_i, z_i, c_i),
\]
where $Y$ includes only observed $y_i$ in the product. Missing prices simply remove the corresponding likelihood factor while still contributing to learning via $x_i$.

A graphical model for a single property is shown in Figure~\ref{fig:graphical-model-miwae}.

\begin{figure}[t]
    \centering
    \includegraphics[width=0.8\textwidth]{miwae_graphical_model.png}
    \caption{Graphical model of the semi-supervised MIWAE with a Student-\emph{t} mixture prior.}
    \label{fig:graphical-model-miwae}
\end{figure}

\subsection*{(iii) Inference}

Direct maximum-likelihood estimation is intractable because the posterior $p(z_i, c_i \mid x_i, y_i)$ is non-Gaussian and multimodal. I therefore use amortized variational inference with an importance-weighted objective (MIWAE):

\paragraph{Variational family.} For each property I introduce an encoder $q_\phi(z_i \mid x_i, y_i)$, parameterized as a diagonal-covariance Gaussian
\[
q_\phi(z_i \mid x_i, y_i) = \mathcal{N}\big(m_\phi(x_i, y_i), \operatorname{diag}(v_\phi(x_i, y_i))\big),
\]
with $m_\phi(\cdot)$ and $\log v_\phi(\cdot)$ parameterized by a neural network. The discrete mixture indicator $c_i$ is marginalized out analytically through the prior; I do not introduce a separate variational distribution over $c_i$.

\paragraph{MIWAE objective.} For partially observed $(x_i, y_i)$ pairs, the training objective aggregates $K_{\text{IW}}$ importance samples from $q_\phi$ and maximizes a Monte Carlo estimate of the log marginal likelihood:
\[
\mathcal{L}_{\text{MIWAE}}(\theta,\phi)
= \sum_{i=1}^n
\mathbb{E}_{z_i^{(1:K_{\text{IW}})} \sim q_\phi}
\left[
\log \left(
\frac{1}{K_{\text{IW}}}
\sum_{k=1}^{K_{\text{IW}}}
\frac{p_\theta(x_i, y_i \mid z_i^{(k)})\, p_\theta(z_i^{(k)})}
     {q_\phi(z_i^{(k)} \mid x_i, y_i)}
\right)
\right].
\]
Here $\theta$ collects all generative parameters: the mixture prior, decoder, and prediction network. In the implementation, the loss is decomposed into a reconstruction term for $x_i$, a supervised log-likelihood term for $y_i$, and a KL regularizer between $q_\phi(z_i \mid x_i, y_i)$ and the mixture prior. A hyperparameter $\alpha_y = 10$ emphasizes the price loss to ensure that the latent variables remain predictive of $y_i$.

\paragraph{Optimization and convergence.} I train the model with Adam on mini-batches, using $K_{\text{IW}}$ importance samples per datapoint and early stopping based on a validation estimate of the MIWAE objective. Convergence diagnostics (train and validation loss per epoch) are reported in Figure~\ref{fig:convergence-loss-miwae}. Across cross-validation folds, the number of epochs to convergence is on the order of \emph{[TODO: insert mean epochs here]}.

\begin{figure}[t]
    \centering
    \includegraphics[width=0.8\textwidth]{miwae_loss_elbo.png}
    \caption{Convergence diagnostics for the SemiSupMIWAE model: training and validation loss versus epoch.}
    \label{fig:convergence-loss-miwae}
\end{figure}

\subsection*{(iv) Data and setup}

\paragraph{Dataset.} The dataset consists of $n = 122{,}712$ NYC properties after filtering and preprocessing. Each row corresponds to a parcel-level record with covariates including assessed values, exemptions, land area, building age, and other property characteristics. I use $D = 43$ input features $x_i$, including:

\begin{itemize}[noitemsep]
    \item Assessed and exempt values: \texttt{assess\_total}, \texttt{assessland}, \texttt{exempt\_total}, etc.
    \item Physical attributes: \texttt{lotarea}, \texttt{yearbuilt}, and related building characteristics.
    \item Additional administrative and categorical encodings (one-hot or continuous) as provided in the original panel.
\end{itemize}

The target $y_i = \log(\texttt{sale\_price}_i)$ is observed for approximately $75\%$ of properties, with the remaining $25\%$ treated as missing and used only through their $x_i$ in the unsupervised MIWAE objective.

\paragraph{Preprocessing.} I standardize all continuous covariates to zero mean and unit variance and log-transform sale prices. Missing covariates are handled through a zero-masking strategy rather than explicit imputation: the encoder and decoder receive binary masks indicating which entries of $x_i$ are observed, and the reconstruction loss is computed only over observed entries. For evaluation, I restrict attention to properties with observed $y_i$.

\paragraph{Model configuration.} The main hyperparameters are:
\begin{itemize}[noitemsep]
    \item Latent dimension: $K = 3$.
    \item Prior: Student-\emph{t} mixture with $K_{\text{prior}} = 5$ components, degrees of freedom fixed at $\nu = 4$.
    \item Encoder/decoder: one hidden layer of width 21 for both $x$-decoder and $z$-encoder.
    \item Prediction network: two hidden layers of widths (8, 4) on top of $z$.
    \item Semi-supervised weight: $\alpha_y = 10.0$ on the price log-likelihood term.
\end{itemize}

I use random 80/20 splits of properties with observed $y_i$ for out-of-sample evaluation when cross-validation predictions are unavailable.

\subsection*{(v) Results}

I report both quantitative predictive metrics and qualitative visualizations of the latent factors.

\subsubsection*{Predictive performance}

For properties with observed sale prices, I evaluate the posterior predictive distribution $p(y_i \mid x_i)$ induced by the trained model. On a random 20\% hold-out set (conditional on $y_i$ being observed), the approximate Gaussian posterior in log-price space yields the following metrics (Table~\ref{tab:global-metrics}):

\begin{itemize}[noitemsep]
    \item Average log posterior predictive $\approx -1.15$ in log-price space.
    \item RMSE in log-price $\approx 1.84$ and MAE in log-price $\approx 0.97$.
    \item In level space, RMSE $\approx 2.46\times 10^7$ and MAE $\approx 2.18\times 10^6$, reflecting the extreme scale and skewness of NYC property prices.
    \item A global MAPE on prices that is numerically dominated by a small number of very low-price transactions and therefore not representative of typical performance.
\end{itemize}

\begin{table}[t]
    \centering
    \begin{tabular}{lcc}
        \toprule
        Metric & Value & Notes \\
        \midrule
        Avg.\ log $p(y \mid x)$ (log-price) & $-1.15$ & Random 20\% hold-out \\
        RMSE (log-price) & $1.84$ & Posterior mean $\hat{y}$ \\
        MAE (log-price) & $0.97$ & Posterior mean $\hat{y}$ \\
        RMSE (price) & $2.46 \times 10^7$ & In dollars \\
        MAE (price) & $2.18 \times 10^6$ & In dollars \\
        % TODO: optionally add robust metrics once computed
        Median $|y - \hat{y}|$ (log) & [TODO] & Robust central error \\
        95th pct.\ $|y - \hat{y}|$ (log) & [TODO] & Tail error \\
        \bottomrule
    \end{tabular}
    \caption{Global posterior predictive metrics for the SemiSupMIWAE model on a random hold-out set. Bracketed entries are placeholders for additional robust metrics (e.g., median absolute error and tail quantiles) once computed.}
    \label{tab:global-metrics}
\end{table}

Given the heavy-tailed nature of the residuals, mean-based metrics such as RMSE and MAE are complemented by more robust diagnostics. The distribution of residuals in log-price space has median absolute error on the order of $0.2$--$0.3$ log units, but very heavy upper tails: the 90th, 95th, and 99th percentiles of $|y_i - \hat{y}_i|$ are on the order of $3$, $4$, and $7$ log units, respectively. A histogram of residuals and a QQ-plot against a Normal reference (Figure~\ref{fig:residuals-log}) highlight substantial kurtosis relative to a Gaussian benchmark.

\begin{figure}[t]
    \centering
    \includegraphics[width=0.49\textwidth]{miwae_residuals_histogram.png}
    \includegraphics[width=0.49\textwidth]{miwae_residuals_qq.png}
    \caption{Residual diagnostics for $\log(\text{sale\_price})$: histogram and QQ-plot versus a Normal reference, illustrating heavy tails.}
    \label{fig:residuals-log}
\end{figure}

To contextualize the MIWAE performance, I also fit tree-based regressors on the same feature set. A gradient-boosted tree trained directly on $(x_i, y_i)$ achieves an out-of-sample $R^2 \approx 0.96$ in log-price space, while a tree trained on the learned latent means $z_i$ attains $R^2 \approx 0.93$. The MIWAE price head itself reaches $R^2 \approx 0.90$ on the same test split. These comparisons suggest that (i) the encoder and prior preserve most of the predictive signal in three latent dimensions, and (ii) the main performance gap relative to the tree upper bound arises from the Gaussian likelihood and small price head rather than the latent representation.

\subsubsection*{Residual structure and unused signal}

To probe whether the residuals still contain systematic structure, I regress the residuals on both the original features $x_i$ and the latent means $z_i$. On a held-out set, a flexible regressor achieves $R^2 \approx 0.50$ for $\text{residual} \mid X$ and $R^2 \approx 0.27$ for $\text{residual} \mid z$, indicating that the model leaves roughly half of the explainable variation in log-price on the table relative to what is available in the raw covariates. This residual structure likely reflects a combination of model misspecification (Gaussian likelihood in log space) and limited capacity in the decoder and price head.

\subsubsection*{Latent structure and qualitative visualization}

Finally, I inspect the learned latent factors. For properties with observed $y_i$, I compute the posterior mean $\mu_z(x_i, y_i)$ and plot the first two dimensions $(z_{i0}, z_{i1})$ colored by sale-price deciles (Figure~\ref{fig:latent-price}).

\begin{figure}[t]
    \centering
    \includegraphics[width=0.8\textwidth]{miwae_latents_by_sale_price.png}
    \caption{Latent space visualization: posterior means $(z_0, z_1)$ colored by sale-price deciles. A smooth gradient along one axis indicates that one latent dimension captures a value/size continuum.}
    \label{fig:latent-price}
\end{figure}

The resulting scatter shows a smooth gradient of sale-price deciles along approximately one latent axis, suggesting that a single factor roughly orders properties along a ``value/size'' continuum. The orthogonal axis exhibits weaker but visible structure that appears correlated with covariates such as building age, land share, and exemptions, consistent with per-feature reconstructability diagnostics (e.g., moderate $R^2$ for reconstructing \texttt{assess\_total}, \texttt{exempt\_total}, \texttt{lotarea}, and \texttt{yearbuilt} from $z$).

When building-class information is available (e.g., small 1--4 family homes, large elevator buildings, mixed-use parcels), I repeat the latent scatter colored by building class (Figure~\ref{fig:latent-building-class}). Clusters in this plot indicate that the latent factors also capture coarse structural differences across property types, even though the model is not explicitly conditioned on building class.

\begin{figure}[t]
    \centering
    \includegraphics[width=0.8\textwidth]{miwae_latents_by_building_class.png}
    \caption{Latent space visualization: posterior means $(z_0, z_1)$ colored by building class (e.g., 1--4 family, large elevator, mixed-use). Distinct clusters suggest that the latent factors encode property-type structure.}
    \label{fig:latent-building-class}
\end{figure}

Overall, the MIWAE model discovers a compact latent representation that organizes properties along interpretable axes and supports competitive predictive performance relative to a flexible tree baseline, while exposing clear heavy-tailed residual structure and leftover signal for future model refinements (e.g., heavier-tailed likelihoods for log prices or richer price heads).

\section*{Plans and Bottlenecks}

\subsection*{Plans for Final Weeks}
In the remaining weeks, I plan to:
\begin{itemize}
    \item Implement heavier-tailed likelihoods (e.g., Student-t) for the price head to better capture extreme valuations.
    \item Conduct a rigorous calibration assessment of the predictive uncertainty.
    \item Explore richer architectures for the price prediction network.
    \item Complete the computation of robust metrics (median absolute error, tail quantiles) to replace the current placeholders.
\end{itemize}

\subsection*{Current Bottlenecks}
The primary bottleneck is the significant ``unused signal'' observed in the residual analysis, where 50\% of the explainable variation remains uncaptured by the latent factors. This suggests potential model misspecification or optimization difficulties with the mixture prior. Additionally, the computational cost of importance sampling limits the number of particles we can use during training, and the heavy-tailed nature of the data continues to pose challenges for standard Gaussian likelihood assumptions.

\section*{Future Work}
Beyond the immediate plans, future work could involve incorporating spatial dependencies explicitly (e.g., via a graph neural network or spatial random effects) and exploring fully unsupervised pre-training on a larger set of unlabeled parcels to improve the latent representation.

\end{document}
